% =====================================================================
%  Module Data Lakehouse -- TP 4 : Batch ou streaming
%  Compilation : pdflatex (2 passes)
% =====================================================================
\documentclass[11pt,a4paper]{article}

\usepackage[utf8]{inputenc}
\usepackage[T1]{fontenc}
\usepackage[french]{babel}
\usepackage{lmodern}
\usepackage[babel=true]{microtype}

\usepackage[a4paper,margin=2.3cm,headheight=14pt]{geometry}
\usepackage{fancyhdr}
\usepackage{titlesec}
\usepackage{needspace}
\usepackage{enumitem}
\usepackage{amssymb}

\usepackage[table]{xcolor}
\usepackage{booktabs}
\usepackage{tabularx}
\usepackage{array}
\usepackage{tikz}
\usetikzlibrary{positioning,arrows.meta,shapes.geometric,babel}
\usepackage{listings}

\usepackage[htt]{hyphenat}
\usepackage[hidelinks]{hyperref}
\usepackage{xurl}
\usepackage{caption}
\usepackage{float}
\usepackage{graphicx}
% Captures d'écran : <nom>.png à la racine du projet (Overleaf) ou dans images/tp4/ ;
% tant que le fichier manque, un cadre indique le nom attendu.
\newcommand{\capture}[3]{%
  \begin{figure}[H]\centering
  \IfFileExists{#1.png}%
    {\includegraphics[width=0.92\textwidth,height=0.42\textheight,keepaspectratio]{#1.png}}%
    {\IfFileExists{images/tp4/#1.png}%
      {\includegraphics[width=0.92\textwidth,height=0.42\textheight,keepaspectratio]{images/tp4/#1.png}}%
      {\fbox{\parbox[c][3.2cm][c]{0.88\textwidth}{\centering\footnotesize
         Capture à ajouter : \texttt{\detokenize{#1}.png}}}}}%
  \caption{#2}\label{#3}
  \end{figure}}
\captionsetup{font=small,labelfont=bf,labelsep=period}

\setlength{\parskip}{0.5em}
\setlength{\parindent}{0pt}
\setlength{\emergencystretch}{3em}

% ---------- Couleurs ----------
\definecolor{lhblue}{HTML}{1F4E79}
\definecolor{lhteal}{HTML}{0F766E}
\definecolor{lhred}{HTML}{B91C1C}
\definecolor{lhviolet}{HTML}{6D28D9}
\definecolor{lhgray}{HTML}{F3F4F6}
\definecolor{lhbronze}{HTML}{A16207}
\definecolor{lhsilver}{HTML}{64748B}
\definecolor{lhgold}{HTML}{CA8A04}
\definecolor{lhfonce}{HTML}{17394D}
\definecolor{lhvert}{HTML}{5FB89D}
% Logos des logiciels (docs/images/logos/, voir LISEZMOI.md) ; sans le fichier, rien ne s'affiche
\newcommand{\logo}[2][1.15em]{\IfFileExists{images/logos/#2.png}{\raisebox{-0.22em}{\includegraphics[height=#1]{images/logos/#2.png}}\hspace{0.3em}}%
  {\IfFileExists{logos/#2.png}{\raisebox{-0.22em}{\includegraphics[height=#1]{logos/#2.png}}\hspace{0.3em}}{}}}
% Pictogrammes de la figure d'architecture (dessins génériques)
\tikzset{
  pics/fichier/.style={code={%
    \fill[white,draw=lhfonce,line width=0.7pt] (-0.32,-0.42) -- (0.18,-0.42) -- (0.32,-0.28) -- (0.32,0.42) -- (-0.32,0.42) -- cycle;
    \draw[lhfonce,line width=0.7pt] (0.18,-0.42) -- (0.18,-0.28) -- (0.32,-0.28);
    \node[font=\sffamily\bfseries\tiny,text=lhfonce] at (0,0.05) {#1};}},
  pics/table/.style={code={%
    \node[cylinder,shape border rotate=90,aspect=0.28,draw=#1!75!black,fill=#1!35,
          minimum width=0.85cm,minimum height=0.95cm,line width=0.6pt] {};}},
  pastille/.style={circle,fill=lhvert,draw=white,line width=1pt,text=white,
                   font=\sffamily\bfseries\scriptsize,inner sep=1.6pt},
  case/.style={draw=lhfonce,fill=white,minimum width=0.4cm,minimum height=0.4cm,inner sep=0pt,
               font=\sffamily\tiny,text=lhfonce},
}


% ---------- Titres, en-têtes ----------
\titleformat{\section}{\needspace{6\baselineskip}\Large\bfseries\color{lhblue}}{\thesection.}{0.6em}{}
\titleformat{\subsection}{\large\bfseries\color{lhblue}}{\thesubsection}{0.6em}{}
\pagestyle{fancy}
\fancyhf{}
\fancyhead[L]{\small\color{lhblue}Module Data Lakehouse}
\fancyhead[R]{\small\color{lhblue}TP 4 -- Batch ou streaming}
\fancyfoot[C]{\small\thepage}

% ---------- Encadrés : texte simple, sans boîte ----------
% Un titre en gras introduit le paragraphe ; seul « Attention » est en rouge.
\newenvironment{resultat}{\par\medskip\noindent\textbf{Résultat attendu :}\ }{\par\medskip}
\newenvironment{astuce}{\par\medskip\noindent\textbf{Astuce :}\ }{\par\medskip}
\newenvironment{attention}[1][]{\par\medskip\noindent\textbf{\textcolor{lhred}{Attention\if\relax\detokenize{#1}\relax\else{} -- #1\fi{} :}}\ }{\par\medskip}

% ---------- Code ----------
\lstset{
  basicstyle=\ttfamily\small, columns=fullflexible, keepspaces=true,
  breaklines=true, showstringspaces=false, upquote=true,
  frame=single, framesep=5pt, xleftmargin=6pt, xrightmargin=6pt,
  backgroundcolor=\color{lhgray}, rulecolor=\color{black!15},
  commentstyle=\color{black!55}\itshape,
  keywordstyle=\color{lhblue}\bfseries, stringstyle=\color{black!70},
  literate=
   {é}{{\'e}}1 {è}{{\`e}}1 {ê}{{\^e}}1 {à}{{\`a}}1 {ç}{{\c c}}1
   {ù}{{\`u}}1 {î}{{\^i}}1 {ô}{{\^o}}1 {É}{{\'E}}1 {Ê}{{\^E}}1
}
\lstdefinestyle{shell}{language=bash,backgroundcolor=\color{black!85},
  basicstyle=\ttfamily\small\color{white},rulecolor=\color{black!85},
  commentstyle=\color{black!35}\itshape,keywordstyle=\color{white}}

\newcommand{\cmd}[1]{\texttt{#1}}
\newcommand{\etape}[1]{\tikz[baseline=(n.base)]\node[circle,fill=lhvert,text=white,
  font=\sffamily\bfseries\scriptsize,inner sep=1.4pt](n){#1};}
\pdfstringdefDisableCommands{\def\cmd#1{#1}}
\hypersetup{pdftitle={TP 4 -- Batch ou streaming},pdfauthor={Anouar Riadsolh}}

\newenvironment{notion}[1][]{\par\medskip\noindent\textbf{\textcolor{lhblue}{Notion clé\if\relax\detokenize{#1}\relax\else{} -- #1\fi{} :}}\ }{\par\medskip}
\newcommand{\nb}[1]{\textcolor{lhblue}{\texttt{#1}}}
% Captures : comment les obtenir (avant la figure), comment les lire (après la figure)
\newenvironment{obtenir}{\par\needspace{12\baselineskip}\smallskip
  \begin{enumerate}[leftmargin=*,nosep,label=\arabic*.]}{\end{enumerate}\par}
\newenvironment{lecture}{\par\noindent\ignorespaces}{\par\medskip}

% =====================================================================
\begin{document}
% =====================================================================

\begin{tikzpicture}[remember picture,overlay]
  \fill[lhblue] (current page.north west) rectangle ([yshift=-5.2cm]current page.north east);
  \node[anchor=north west,text=white,font=\sffamily] at ([xshift=2.3cm,yshift=-1.3cm]current page.north west)
    {Module Data Lakehouse};
  \node[anchor=north west,text=white,font=\sffamily\bfseries\fontsize{26}{30}\selectfont] at ([xshift=2.3cm,yshift=-2.2cm]current page.north west)
    {TP 4 -- Batch ou streaming};
  \node[anchor=north west,text=white!85!lhblue,font=\sffamily\large] at ([xshift=2.3cm,yshift=-3.6cm]current page.north west)
    {Des événements en continu : Kafka, Spark Structured Streaming et Iceberg};
  \node[anchor=north east,text=white!85!lhblue,font=\sffamily] at ([xshift=-2.3cm,yshift=-4.25cm]current page.north east)
    {Pr.~Anouar Riadsolh};
\end{tikzpicture}
\thispagestyle{empty}
\vspace*{4.2cm}

{\large\bfseries\color{lhblue}En bref}\par\smallskip
\begin{tabularx}{\linewidth}{@{}>{\bfseries}l X@{}}
\toprule
Objectif & Recevoir le suivi des colis sous forme d'événements, en continu, avec Kafka, l'écrire dans
           le Lakehouse sans perte ni doublon avec Spark Structured Streaming, traiter les événements
           en retard, puis décider, mesures à l'appui, quand le streaming vaut son coût. \\
Durée    & 2 séances de 3 h (+ environ 2 h de travail personnel). \\
Prérequis & TP 1 à 3 terminés. Images du cours mises à jour et profil \cmd{streaming}
            (section~\ref{sec:deroule}). \\
À rendre & Un compte rendu et les cinq notebooks exécutés (section~\ref{sec:livrables}). \\
\bottomrule
\end{tabularx}

% =====================================================================
\section{Le scénario}
% =====================================================================

Au TP 2, les transporteurs envoyaient le suivi des colis une fois par mois : un message par colis,
avec toutes ses étapes, une fois le colis livré. Ils proposent maintenant d'envoyer chaque
\textbf{événement} au moment où il se produit : colis pris en charge, en transit, en livraison,
livré, ou incident. Trois équipes sont intéressées :

\begin{enumerate}[leftmargin=*,nosep]
  \item \textbf{Le service client} veut savoir où en est chaque colis et être prévenu d'un incident
        dans la journée, pas un mois plus tard.
  \item \textbf{La direction} veut toujours son bilan mensuel des incidents par transporteur, exact.
  \item \textbf{L'équipe data} doit stocker ces événements dans le Lakehouse sans en perdre ni en
        doubler, même si un serveur tombe en panne.
\end{enumerate}

Les données n'ont plus de fin : c'est un \textbf{flux}. Ce TP termine le module. Il ajoute à
l'architecture construite aux TP 1 à 3 la dimension temps réel, et la question qui va avec : de
quel délai le métier a-t-il vraiment besoin ?

% =====================================================================
\section{Les données}
% =====================================================================

Le flux est reconstitué à partir des messages du TP 2 (\cmd{data/olist/suivi/}). Le script
\cmd{scripts/rejouer\_suivi.py} découpe chaque message en événements, un par étape ou par incident,
et les envoie dans l'ordre où ils se sont produits, en \textbf{temps accéléré} : par défaut, une
journée Olist par seconde. Il écrit dans Kafka, ou dans un dossier avec l'option \cmd{-{}-dossier}.

\begin{lstlisting}
{"evenement_id": "CL5A9F638CAE9F-E1", "colis_id": "CL5A9F638CAE9F",
 "order_id": "5a9f638cae9fc9cd16c2f1e748da4c77", "transporteur": "Correios",
 "categorie": "etape", "code": "EN_TRANSIT", "horodatage": "2018-03-01 00:00:23",
 "ville": "Cajamar", "etat": "SP"}
\end{lstlisting}

\begin{center}
\small
\renewcommand{\arraystretch}{1.25}
\begin{tabularx}{\textwidth}{@{}l X@{}}
\toprule
\textbf{Champ} & \textbf{Contenu} \\
\midrule
\cmd{evenement\_id} & identifiant de l'événement : le colis, puis \cmd{E} et le rang de l'étape, ou \cmd{I} et le rang de l'incident \\
\cmd{categorie}, \cmd{code} & \cmd{etape} (\cmd{PRIS\_EN\_CHARGE}, \cmd{EN\_TRANSIT}, \cmd{EN\_LIVRAISON}, \cmd{LIVRE}) ou \cmd{incident} (\cmd{DESTINATAIRE\_ABSENT}, \cmd{COLIS\_ENDOMMAGE}\dots) \\
\cmd{horodatage} & le moment où l'événement s'est produit \\
\cmd{transporteur}, \cmd{ville}, \cmd{etat} & qui transporte le colis, et où il se trouve \\
\bottomrule
\end{tabularx}
\end{center}

Le mois de mars 2018 compte 34\,951 événements. Comme dans un vrai flux, le rejeu garde trois
défauts, et le tirage est fixé : deux rejeux envoient exactement les mêmes événements.
\begin{itemize}[leftmargin=*,nosep]
  \item \textbf{des messages tronqués} : quelques-uns par mois, au JSON invalide ;
  \item \textbf{des doublons} : un message renvoyé par le transporteur redonne les mêmes événements, avec le même \cmd{evenement\_id} ;
  \item \textbf{des retards} : 5~\% des événements arrivent de 2 à 48 heures après s'être produits.
\end{itemize}

\begin{notion}[Événement et état]
un événement décrit ce qui \emph{s'est passé} (\og le colis est
arrivé en transit à Cajamar à 0 h 00 \fg{}) : on ne le modifie jamais, on en ajoute. L'état décrit
ce qui \emph{est vrai maintenant} (\og le colis est en transit \fg{}) : il change à chaque
événement. Un flux transporte des événements, et c'est au traitement d'en déduire les états.
\end{notion}

% =====================================================================
\section{L'architecture du flux}
% =====================================================================

La figure~\ref{fig:archi} montre le chemin d'un événement, du transporteur jusqu'aux équipes métier.
Elle reprend le Lakehouse des TP précédents, et y ajoute \textbf{Kafka}, qui transporte le flux, et
\textbf{Spark Structured Streaming}, qui le traite par micro-lots.

\begin{figure}[H]
\centering
\resizebox{\textwidth}{!}{%
\begin{tikzpicture}[
  font=\sffamily\small, >=Stealth,
  barre/.style={fill=lhfonce,text=white,font=\sffamily\bfseries\small,minimum height=0.6cm,inner sep=0pt},
  cadre/.style={draw=lhfonce,line width=0.8pt},
  bloc/.style={draw=lhfonce!60,line width=0.6pt,fill=white},
  et/.style={font=\sffamily\scriptsize,text=lhfonce,align=center},
  flux/.style={->,very thick,lhfonce}]

% ===== Sources =====
\draw[cadre] (0,0) rectangle (3.2,-10.2);
\node[barre,minimum width=3.2cm,anchor=north west] at (0,0) {Sources};
\pic at (1.6,-1.55) {fichier=\{ \}};
\node[et] at (1.6,-2.4) {Suivi des colis\\fichiers JSON};
\draw[->,thick,lhfonce] (1.6,-2.85) -- (1.6,-3.15);
\node[draw=lhfonce,line width=0.7pt,fill=lhgray,rounded corners=2pt,font=\sffamily\scriptsize,align=center,
      minimum width=2.4cm] at (1.6,-3.6) {\texttt{rejouer\_suivi.py}\\producteur};
\node[et] at (1.6,-4.75) {un événement\\par étape, en\\temps accéléré};

% ===== Kafka =====
\draw[cadre] (3.9,0) rectangle (7.4,-10.2);
\node[barre,minimum width=3.5cm,anchor=north west] at (3.9,0) {\logo[1.2em]{kafka_blanc}Kafka};
\node[et,font=\sffamily\bfseries\scriptsize] (topic) at (5.65,-1.2) {topic \texttt{suivi\_colis}};
\node[pastille,left=1pt] at (topic.west) {2};
% pas de \foreach dans \resizebox : babel français le casse
\node[et] at (4.3,-1.9) {P0};
\node[case] at (4.75,-1.9) {0}; \node[case] at (5.15,-1.9) {1}; \node[case] at (5.55,-1.9) {2}; \node[case] at (5.95,-1.9) {3}; \node[case] at (6.35,-1.9) {4}; \node[case] at (6.75,-1.9) {5};
\node[et] at (4.3,-2.6) {P1};
\node[case] at (4.75,-2.6) {0}; \node[case] at (5.15,-2.6) {1}; \node[case] at (5.55,-2.6) {2}; \node[case] at (5.95,-2.6) {3}; \node[case] at (6.35,-2.6) {4}; \node[case] at (6.75,-2.6) {5};
\node[et] at (4.3,-3.3) {P2};
\node[case] at (4.75,-3.3) {0}; \node[case] at (5.15,-3.3) {1}; \node[case] at (5.55,-3.3) {2}; \node[case] at (5.95,-3.3) {3}; \node[case] at (6.35,-3.3) {4}; \node[case] at (6.75,-3.3) {5};
\node[et] at (5.65,-4.15) {3 partitions,\\clé \texttt{colis\_id},\\offsets 0, 1, 2\dots};
\node[et] at (5.65,-5.6) {les messages restent\\7 jours après lecture};
\draw[bloc] (4.3,-7.4) rectangle (7.0,-9.6);
\node[et,font=\sffamily\bfseries\scriptsize] at (5.65,-7.85) {Kafka UI};
\node[et] at (5.65,-8.7) {topics, messages,\\retard des groupes\\\texttt{localhost:8090}};

% ===== Plateforme Lakehouse =====
\draw[draw=lhblue,line width=2.4pt] (8.1,0) rectangle (18.6,-10.2);
\node[fill=lhblue,text=white,font=\sffamily\bfseries,minimum width=10.5cm,minimum height=0.7cm,anchor=north west]
  at (8.1,0) {Plateforme Lakehouse};
\node[barre,minimum width=9.6cm] (spark) at (13.35,-1.25)
  {\logo[1.1em]{spark}Spark Structured Streaming \quad \logo[1.1em]{airflow}déclenchement};
\node[pastille] at (spark.north west) {7};

% Tables par couche
\draw[bloc] (8.55,-2.2) rectangle (18.15,-6.15);
\node[et,font=\sffamily\bfseries\scriptsize] at (13.35,-2.55) {Tables par couche};
\pic at (9.7,-3.55)  {table=lhbronze};  \node[et] at (9.7,-4.45)  {bronze\\événements bruts};
\pic at (12.7,-3.55) {table=lhsilver};  \node[et] at (12.7,-4.45) {silver\\sans doublons};
\pic at (16.6,-3.55) {table=lhgold};    \node[et] at (16.6,-4.45) {gold\\incidents par jour};
\pic at (12.7,-5.55) {table=gray};      \node[et,anchor=west] at (13.2,-5.55) {quarantaine};
\pic at (16.6,-5.55) {table=lhsilver};  \node[et,anchor=west] at (17.1,-5.55) {état\\des colis};
\pic at (9.7,-5.55) {fichier=ckpt};     \node[et,anchor=west] at (10.1,-5.55) {checkpoint};
\draw[flux] (10.25,-3.55) -- (12.15,-3.55) node[pastille,midway]{4};
\draw[flux] (13.25,-3.55) -- (16.05,-3.55) node[pastille,midway]{5};
\draw[->,thick,lhfonce] (12.7,-4.8) -- (12.7,-5.0);
\draw[flux] (13.25,-3.9) -- (16.05,-5.3) node[pastille,pos=0.55]{6};

% Socle : catalogue et stockage ouvert
\node[fill=lhblue,text=white,font=\sffamily\bfseries\small,minimum width=9.6cm,minimum height=0.6cm]
  at (13.35,-6.85) {Lakehouse ouvert};
\draw[bloc] (8.55,-7.3) rectangle (11.95,-9.75);
\node[barre,minimum width=3.4cm,anchor=north west] at (8.55,-7.3) {Catalogue};
\node[et] at (10.25,-8.55) {\logo[2.0em]{iceberg}\\[2pt]Iceberg REST};
\draw[bloc] (12.35,-7.3) rectangle (18.15,-9.75);
\node[barre,minimum width=5.8cm,anchor=north west] at (12.35,-7.3) {Stockage ouvert (S3)};
\node[et] at (13.6,-8.55) {\logo[2.0em]{rustfs}\\[2pt]RustFS};
\node[et] at (15.25,-8.55) {\logo[2.0em]{iceberg}\\[2pt]Iceberg};
\node[et] at (16.9,-8.55) {\logo[2.0em]{parquet}\\[2pt]Parquet};

% ===== Usages =====
\draw[cadre] (19.3,0) rectangle (22.7,-4.8);
\node at (21.0,-1.5) {\logo[2em]{trino}\logo[1.8em]{superset}};
\node[et] at (21.0,-2.6) {Service client\\état des colis,\\incidents du jour};
\draw[cadre] (19.3,-5.4) rectangle (22.7,-10.2);
\node at (21.0,-6.9) {\logo[1.9em]{spark}};
\node[et] at (21.0,-8.0) {Direction\\bilan mensuel\\en batch};

% ===== Infrastructure =====
\node[fill=lhfonce,text=white,minimum width=22.7cm,minimum height=0.75cm,anchor=north west] at (0,-10.6)
  {\logo[1.3em]{docker}\textbf{Infrastructure : conteneurs Docker}};

% ===== Flux =====
\draw[flux] (2.8,-3.6) -- (3.9,-3.6) node[pastille,pos=0.75]{1};
\draw[flux] (7.0,-2.6) -- (9.15,-3.3) node[pastille,pos=0.45]{3};
\draw[flux] (18.6,-2.4) -- (19.3,-2.4) node[pastille,midway,above=3pt]{8};
\draw[flux] (18.6,-7.8) -- (19.3,-7.8);
\end{tikzpicture}}
\caption{L'architecture du flux de suivi des colis. Tous les services tournent dans des conteneurs Docker.}
\label{fig:archi}
\end{figure}

Les numéros renvoient aux étapes suivantes.

\begin{description}[leftmargin=1.4em,labelindent=0pt,itemsep=3pt,font=\normalfont]
\item[\etape{1}] \textbf{Production} (\nb{01\_du\_fichier\_au\_flux}, \nb{02\_kafka}) : le script de
  rejeu envoie chaque événement au moment où il se produit, d'abord dans un dossier, puis dans Kafka.
\item[\etape{2}] \textbf{Transport} (\nb{02\_kafka}) : Kafka range les messages du topic
  \cmd{suivi\_colis} dans trois partitions, selon leur clé \cmd{colis\_id}, et les garde après lecture.
\item[\etape{3}] \textbf{Bronze} (\nb{03\_bronze\_streaming}) : Spark lit le topic en continu et écrit
  chaque message dans \cmd{bronze.suivi\_evenements}. Le checkpoint retient les offsets traités.
\item[\etape{4}] \textbf{Silver} (\nb{04\_temps\_fenetres}) : les messages illisibles partent en
  quarantaine, et les doublons sont écartés grâce au \emph{watermark}.
\item[\etape{5}] \textbf{Gold} (\nb{04\_temps\_fenetres}) : les incidents sont comptés par jour et par
  transporteur, dans des fenêtres de temps.
\item[\etape{6}] \textbf{État des colis} (\nb{04\_temps\_fenetres}) : chaque micro-lot met à jour la
  dernière étape de chaque colis, par un \cmd{MERGE}.
\item[\etape{7}] \textbf{Déclenchement} (\nb{05\_batch\_ou\_streaming}) : la même requête tourne en
  continu, ou à heure fixe, lancée par Airflow.
\item[\etape{8}] \textbf{Usages} : le service client lit l'état des colis et les incidents du jour
  avec Trino et Superset, et la direction garde son bilan mensuel, calculé en batch.
\end{description}

% =====================================================================
\section{Déroulé}\label{sec:deroule}
% =====================================================================

\textbf{Mettre à jour l'environnement, une fois, au début de la séance 1.} L'image JupyterLab
contient maintenant le connecteur Kafka, et un nouveau profil \cmd{streaming} démarre Kafka et son
interface web. Dans le terminal Ubuntu, depuis le dossier du dépôt :
\begin{lstlisting}[style=shell]
git pull
cp .env.example .env          # ajoute les versions de Kafka
docker compose --profile full --profile orchestration --profile observabilite stop
docker compose pull jupyter
docker compose --profile streaming pull kafka kafka-ui
docker compose up -d
docker compose --profile streaming up -d kafka kafka-ui
\end{lstlisting}
Si vous aviez modifié \cmd{.env} (machine de 8 Go), refaites la même modification. Puis, dans un
terminal JupyterLab, copiez les énoncés :
\begin{lstlisting}[style=shell]
cp -r notebooks/tp4 notebooks/mon_travail/
\end{lstlisting}

\begin{center}
\small
\renewcommand{\arraystretch}{1.25}
\begin{tabularx}{\textwidth}{@{}l X l@{}}
\toprule
\textbf{Partie} & \textbf{Contenu} & \textbf{Services nécessaires} \\
\midrule
1 & Du fichier au flux (notebook \nb{01\_du\_fichier\_au\_flux}) & de base \\
2 & Transporter le flux avec Kafka (\nb{02\_kafka}) & \cmd{-{}-profile streaming} \\
3 & Écrire le flux dans Bronze (\nb{03\_bronze\_streaming}) & \cmd{-{}-profile streaming} \\
4 & Le temps dans un flux (\nb{04\_temps\_fenetres}) & \cmd{-{}-profile streaming} \\
5 & Batch ou streaming ? (\nb{05\_batch\_ou\_streaming}) & \cmd{-{}-profile streaming} \\
\bottomrule
\end{tabularx}
\end{center}

\begin{attention}[Mémoire]
Kafka et Kafka UI occupent environ 700~Mo. Arrêtez Airflow, Trino, Superset et la supervision
pendant ce TP (première commande ci-dessus), et un seul notebook à la fois : avant d'ouvrir le
suivant, arrêtez le noyau du précédent (\emph{Kernel > Shut Down Kernel}). Le message
\cmd{SparkUI could not bind on port 4040} signale qu'un autre noyau tourne encore.
\end{attention}

% =====================================================================
\section{Partie par partie}
% =====================================================================

\subsection*{Partie 1 -- Du fichier au flux \hfill \nb{01\_du\_fichier\_au\_flux}}

Le rejeu dépose un fichier JSON par jour dans \cmd{landing/flux/suivi/}, toutes les deux secondes.
Vous lisez ce dossier \emph{en continu} et vous comptez les événements par code : la requête se met
à jour seule à chaque nouveau fichier.

\begin{notion}[Spark Structured Streaming]
Spark traite un flux comme une \textbf{table qui grandit sans
fin}. On écrit la requête comme en batch, avec \cmd{readStream} au lieu de \cmd{read} et
\cmd{writeStream} au lieu de \cmd{write}. Spark la relance ensuite à chaque \textbf{micro-lot}, sur
les seules données nouvelles, et met le résultat à jour. Pour une agrégation, il garde d'un
micro-lot à l'autre un \textbf{état}, par exemple le compte courant de chaque code.
\end{notion}

\begin{center}
\small
\renewcommand{\arraystretch}{1.25}
\begin{tabularx}{\textwidth}{@{}l X l@{}}
\toprule
\textbf{Mode de sortie} & \textbf{Ce qui est écrit à chaque micro-lot} & \textbf{Avec une agrégation} \\
\midrule
\cmd{append}   & les nouvelles lignes, définitives & avec un \emph{watermark} seulement \\
\cmd{update}   & les lignes modifiées & oui \\
\cmd{complete} & tout le résultat & oui \\
\bottomrule
\end{tabularx}
\end{center}

\begin{resultat}
les 11\,232 événements du 1er au 10 mars 2018 sont lus en 10 micro-lots, un par fichier, sans en
oublier ni en lire deux fois. Sans schéma, Spark refuse de démarrer, et le mode \cmd{append} est
refusé sur une agrégation sans \emph{watermark}.
\end{resultat}

\begin{obtenir}
  \item Dans \nb{01\_du\_fichier\_au\_flux}, exécutez les cellules 1 à 5.
  \item Exécutez la cellule 6 : elle affiche le résultat de la requête toutes les 5 secondes, pendant que les fichiers arrivent.
  \item Capturez la sortie de la cellule 6, des tableaux \og après 10 s \fg{} à \og après 20 s \fg{}.
\end{obtenir}
\capture{tp4_01_micro_lots}{Partie 1 : le nombre d'événements par code, mis à jour à chaque micro-lot.}{fig:res01}
\begin{lecture}
Le même tableau est affiché trois fois, à 5 secondes d'intervalle, et ses compteurs grossissent :
\cmd{EN\_TRANSIT} passe de 1\,348 à 2\,116 puis à 2\,879 événements, à mesure que les micro-lots 2 à 7
lisent les fichiers des jours suivants. Personne n'a relancé la requête : Spark l'a fait à chaque
nouveau fichier, en ne lisant que lui, et en ajoutant son résultat à l'état gardé en mémoire. À la
fin du rejeu, le total lu est de 11\,232, exactement le nombre d'événements envoyés.
\end{lecture}


\subsection*{Partie 2 -- Transporter le flux avec Kafka \hfill \nb{02\_kafka}}

Vous créez le topic \cmd{suivi\_colis}, vous y envoyez une semaine d'événements, puis vous les lisez
avec deux groupes de consommateurs (figure~\ref{fig:kafka}).

\begin{figure}[H]
\centering
\begin{tikzpicture}[font=\sffamily\footnotesize, >=Stealth,
  et/.style={font=\sffamily\scriptsize,text=lhfonce,align=center}]
\node[draw=lhfonce,fill=lhgray,rounded corners=2pt,align=center,minimum height=1.6cm,text width=2.3cm]
  (prod) at (0,-1.0) {\textbf{producteur}\\clé \texttt{colis\_id}\\$\rightarrow$ partition};
\foreach \p/\y/\n in {0/0/6,1/-1.0/5,2/-2.0/4} {
  \node[et] at (2.3,\y) {partition \p};
  \foreach \o in {0,...,\n} \node[case,minimum width=0.55cm,minimum height=0.55cm] at (3.3+\o*0.55,\y) {\o};
  \draw[->,thick,lhfonce] (prod.east) -- (2.3-0.65,\y);
}
\node[et,font=\sffamily\bfseries\scriptsize] at (5.0,0.75) {topic \texttt{suivi\_colis}};
\draw[->,thick,lhteal] (8.6,0.55) -- (5.0,0.3);
\node[et,text=lhteal,anchor=west] at (8.6,0.55) {groupe \texttt{tp4\_lecteurs} : position 2};
\draw[->,thick,lhviolet] (8.6,-2.55) -- (5.5,-2.3);
\node[et,text=lhviolet,anchor=west] at (8.6,-2.55) {groupe \texttt{tp4\_tableau\_de\_bord} : à jour};
\end{tikzpicture}
\caption{Un topic de trois partitions. Le producteur choisit la partition d'après la clé, et chaque
groupe de consommateurs garde sa propre position (son offset) dans chaque partition.}
\label{fig:kafka}
\end{figure}

\begin{notion}[Le journal de Kafka]
un topic est un \textbf{journal} : les messages y sont ajoutés à la
fin et numérotés (l'\textbf{offset}), puis gardés pendant la durée de rétention, qu'ils aient été
lus ou non. L'ordre n'est garanti que dans une \textbf{partition}, et la \textbf{clé} décide de la
partition : tous les événements d'un colis restent ainsi dans l'ordre de leur envoi. Chaque
\textbf{groupe} de consommateurs mémorise sa position et lit à son rythme. Son \textbf{retard}
(\emph{lag}) est le nombre de messages qui lui restent à lire.
\end{notion}

\begin{resultat}
la première semaine de mars compte 7\,869 messages, répartis en 2\,718, 2\,639 et 2\,512 entre les
trois partitions. Aucun des 4\,319 colis n'est réparti sur deux partitions. Après la lecture, le
groupe \cmd{tp4\_lecteurs} a encore 7\,859 messages de retard, et \cmd{tp4\_tableau\_de\_bord} n'en a
aucun.
\end{resultat}

\begin{obtenir}
  \item Dans \nb{02\_kafka}, exécutez les cellules 1 à 3, qui créent le topic et y envoient la première semaine de mars.
  \item Ouvrez Kafka UI (\url{http://localhost:8090}), menu \emph{Topics}, puis le topic \cmd{suivi\_colis} et l'onglet \emph{Messages}.
  \item Capturez la liste des messages, triée du plus récent au plus ancien (\emph{Newest}).
\end{obtenir}
\capture{tp4_02_messages}{Partie 2 : les messages du topic \texttt{suivi\_colis} dans Kafka UI.}{fig:res02a}
\begin{lecture}
Chaque ligne est un message : son offset, sa partition, son heure d'arrivée dans Kafka, sa clé et
sa valeur, le JSON de l'événement. Les derniers offsets sont 2\,717 pour la partition 0 et 2\,511
pour la partition 2 : ces partitions contiennent 2\,718 et 2\,512 messages, comme l'indique la
cellule 4. La clé est toujours un \cmd{colis\_id}, et la valeur commence par l'\cmd{evenement\_id} de
ce colis. Tous les messages ont la même heure d'arrivée, à la seconde près : le rejeu a été lancé
sans pause. C'est la différence entre l'heure d'arrivée, que Kafka connaît, et l'heure de
l'événement, cachée dans la valeur, qui sera au cœur de la partie 4.
\end{lecture}

\begin{obtenir}
  \item Dans \nb{02\_kafka}, exécutez les cellules 4 à 7.
  \item Dans Kafka UI, ouvrez le menu \emph{Consumers}.
  \item Capturez la liste des groupes, avec leur retard (\emph{Consumer Lag}).
\end{obtenir}
\capture{tp4_02_consumers}{Partie 2 : les groupes de consommateurs et leur retard dans Kafka UI.}{fig:res02b}
\begin{lecture}
Le groupe \cmd{tp4\_tableau\_de\_bord} a tout lu : son retard est nul. Le groupe
\cmd{tp4\_lecteurs} n'a lu que les 10 messages de l'étape 3 : son retard est de 7\,859 messages, soit
tout le topic moins ces 10 messages. Les deux groupes lisent le même topic sans se gêner, et les
messages lus par l'un restent disponibles pour l'autre. En production, un retard qui grandit est le
premier signe qu'un consommateur ne suit plus le rythme du flux.
\end{lecture}


\subsection*{Partie 3 -- Écrire le flux dans Bronze \hfill \nb{03\_bronze\_streaming}}

Spark lit le topic en continu et écrit chaque message dans \cmd{bronze.suivi\_evenements}, avec sa
partition, son offset et son heure d'arrivée. En plein rejeu, vous arrêtez la requête, comme si le
serveur tombait en panne, puis vous la relancez.

\begin{notion}[Le checkpoint et la garantie \emph{exactly-once}]
à chaque micro-lot, Spark note dans
le \textbf{checkpoint} (un dossier dans RustFS) les offsets qu'il va lire, puis ceux qu'il a écrits.
Iceberg n'enregistre un micro-lot que par un \textbf{commit atomique} : un micro-lot est écrit en
entier, ou pas du tout. Après une panne, Spark relit le checkpoint et reprend au premier offset non
écrit. Chaque message est donc écrit \textbf{une fois et une seule}. Un checkpoint appartient à
\emph{une} requête : le supprimer fait tout relire, et le partager entre deux requêtes les fausse.
\end{notion}

\begin{resultat}
le mois de mars, soit 34\,951 messages dont 9 tronqués, arrive en Bronze. Après la panne, le premier
micro-lot rattrape d'un coup plus de 15\,000 messages, et la table compte autant de lignes que de
messages envoyés, toutes à des positions Kafka différentes. Un mois de flux laisse déjà 7 ou 8
snapshots et une vingtaine de petits fichiers.
\end{resultat}

\begin{obtenir}
  \item Dans \nb{03\_bronze\_streaming}, exécutez les cellules 1 à 4 sans attendre : le rejeu dure 30 secondes.
  \item Exécutez les cellules 5 à 7, qui simulent la panne, relancent la requête, puis vérifient la table.
  \item Capturez les sorties des cellules 5 à 7.
\end{obtenir}
\capture{tp4_03_panne}{Partie 3 : la panne, la reprise et la vérification de la table Bronze.}{fig:res03}
\begin{lecture}
Au moment de la panne, Bronze contient 5\,579 lignes, et toujours 5\,579 dix secondes plus tard :
plus rien n'est écrit, mais le rejeu continue d'envoyer dans Kafka, où les messages s'accumulent. À
la reprise, le micro-lot 3 lit d'un coup 15\,889 messages, contre 1\,000 à 5\,000 pour un micro-lot
normal : Spark a repris aux offsets notés dans le checkpoint. Au final, 34\,951 messages ont été
envoyés, la table compte 34\,951 lignes et 34\,951 positions distinctes : aucune perte, aucun
doublon.
\end{lecture}


\subsection*{Partie 4 -- Le temps dans un flux \hfill \nb{04\_temps\_fenetres}}

Vous mesurez le désordre du flux, puis vous construisez la couche Silver (quarantaine, doublons),
un indicateur par fenêtre de temps en Gold, et l'état courant de chaque colis.

\begin{figure}[H]
\centering
\begin{tikzpicture}[font=\sffamily\footnotesize, >=Stealth,
  et/.style={font=\sffamily\scriptsize,text=lhfonce,align=center}]
\draw[->,thick,lhfonce] (0,0) -- (13,0) node[right,et]{heure de\\l'événement};
\fill[lhred!12] (0,-0.05) rectangle (6.0,0.9);
\fill[lhvert!20] (6.0,-0.05) rectangle (10.0,0.9);
\draw[lhred,very thick] (6.0,-0.3) -- (6.0,1.1);
\node[et,text=lhred,anchor=north] at (6.0,-0.3) {\textbf{watermark}\\= plus récent vu $-$ attente};
\draw[lhfonce,very thick] (10.0,-0.3) -- (10.0,1.1);
\node[et,anchor=north] at (10.0,-0.3) {\textbf{événement le plus récent}\\\textbf{déjà vu}};
\draw[<->,lhfonce] (6.05,1.9) -- node[et,above]{attente (par exemple 12 h)} (9.95,1.9);
\foreach \x in {6.6,7.4,8.3,9.1,9.7} \fill[lhfonce] (\x,0.4) circle (0.09);
\fill[lhteal] (7.9,0.65) circle (0.12);
\node[et,text=lhteal,anchor=south] at (7.9,0.8) {en retard,\\accepté};
\fill[lhred] (3.6,0.4) circle (0.12);
\node[et,text=lhred] at (3.6,0.95) {trop en retard : ignoré};
\node[et] at (2.4,-0.55) {journées closes : écrites,\\elles ne changeront plus};
\end{tikzpicture}
\caption{Le \emph{watermark} : Spark attend les retardataires pendant une durée fixée. Un événement
plus ancien que le watermark arrive trop tard : sa journée est déjà close.}
\label{fig:watermark}
\end{figure}

\begin{notion}[Heure de l'événement et watermark]
l'\textbf{heure de l'événement} est le moment où
le colis a changé d'état (\cmd{horodatage}). L'\textbf{heure de traitement} est le moment où le
message arrive dans Kafka ou dans Spark. Les indicateurs se calculent sur la première, mais les
événements n'arrivent pas toujours dans l'ordre. Le \textbf{watermark}
(\cmd{withWatermark("horodatage\_ts", "12 hours")}) dit combien de temps attendre les retardataires
(figure~\ref{fig:watermark}). Il borne aussi l'\textbf{état} de Spark : au-delà, les fenêtres sont
closes et les identifiants déjà vus sont oubliés.
\end{notion}

\begin{notion}[Mettre à jour une table depuis un flux]
le mode \cmd{append} ne sait qu'ajouter des
lignes. Pour tenir à jour une ligne par colis, on confie chaque micro-lot à une fonction Python
(\cmd{foreachBatch}), qui exécute un \cmd{MERGE} Iceberg. La condition
\cmd{s.horodatage\_ts > t.horodatage\_ts} empêche un événement en retard d'écraser un état plus récent.
\end{notion}

\begin{resultat}
1\,806 événements sur 34\,942 lisibles (5,2~\%) arrivent après un événement plus récent qu'eux, dont
991 avec plus de 24 heures de retard. Silver garde 34\,755 événements, 9 messages vont en quarantaine
et 187 doublons sont écartés. L'état des colis compte une ligne pour chacun des 10\,118 colis,
dont 6\,835 livrés.
\end{resultat}

\begin{obtenir}
  \item Dans \nb{04\_temps\_fenetres}, exécutez les cellules 1 à 5.
  \item Exécutez la cellule 6, qui compare les incidents par jour calculés en flux, avec 12 heures d'attente, au calcul batch, puis la cellule 7, qui refait la comparaison avec 2 jours d'attente.
  \item Capturez les sorties des cellules 6 et 7.
\end{obtenir}
\capture{tp4_04_watermark}{Partie 4 : les incidents par jour, en flux avec 12 heures puis 2 jours d'attente, comparés au batch.}{fig:res04}
\begin{lecture}
Avec 12 heures d'attente, 30 journées sur 31 sont closes, et le flux compte 1\,075 incidents contre
1\,084 pour le batch : 9 incidents (0,8~\%) sont arrivés après la clôture de leur journée, et ne sont
pas comptés. Avec 2 jours d'attente, l'écart tombe à 0, mais seules 28 journées sont closes : le
chiffre de chaque journée arrive deux jours plus tard. Le watermark règle donc un compromis entre
\textbf{fraîcheur} et \textbf{exactitude}, et sa durée se choisit d'après le retard réellement
observé dans le flux.
\end{lecture}


\subsection*{Partie 5 -- Batch ou streaming ? \hfill \nb{05\_batch\_ou\_streaming}}

Vous calculez le même indicateur, les incidents par transporteur, de trois façons, avec la même
fonction Python, puis vous comparez leur délai et leur coût (figure~\ref{fig:modes}).

\begin{figure}[H]
\centering
\begin{tikzpicture}[font=\sffamily\footnotesize, >=Stealth,
  et/.style={font=\sffamily\scriptsize,text=lhfonce,align=left},
  bloc/.style={draw=lhfonce,fill=#1,minimum height=0.42cm,inner sep=0pt}]
\foreach \x in {0.3,0.9,1.2,2.0,2.6,3.1,3.9,4.4,5.2,5.6,6.3,7.0,7.4,8.1,8.8,9.2,9.9,10.5,11.1,11.6}
  \fill[lhfonce] (2.6+\x,1.0) circle (0.06);
\node[et,anchor=east] at (2.4,1.0) {événements};
\node[et,anchor=east] at (2.4,0.3) {batch nocturne};
\node[bloc=lhgold!60,minimum width=1.2cm] at (13.6,0.3) {};
\node[et,anchor=east] at (2.4,-0.4) {flux continu};
\foreach \x in {0,...,23} \node[bloc=lhvert!70,minimum width=0.28cm] at (2.75+\x*0.5,-0.4) {};
\node[et,anchor=east] at (2.4,-1.1) {flux déclenché};
\foreach \x in {5.6,8.6,11.6,14.6} \node[bloc=lhteal!45,minimum width=0.55cm] at (\x,-1.1) {};
\draw[->,thick,lhfonce] (2.6,-1.6) -- (14.9,-1.6) node[et,anchor=north east]{temps};
\end{tikzpicture}
\caption{Trois façons de traiter le même flux : tout relire une fois par nuit, un micro-lot toutes
les quelques secondes, ou les seuls nouveaux messages, à heure fixe.}
\label{fig:modes}
\end{figure}

\begin{notion}[Le déclencheur, ou \emph{trigger}]
le même code Structured Streaming tourne de deux
façons, selon une seule option. Avec \cmd{trigger(processingTime="5 seconds")}, la requête ne
s'arrête jamais et lance un micro-lot toutes les 5 secondes. Avec \cmd{trigger(availableNow=True)},
elle traite tout ce qui est arrivé depuis son dernier passage, puis s'arrête : c'est un
\textbf{batch incrémental}, qu'Airflow peut lancer chaque heure. Le checkpoint garde les offsets et
l'état d'un passage à l'autre.
\end{notion}

\begin{resultat}
le flux continu écrit le dernier incident dans la table quelques secondes après son envoi. Le flux
déclenché ne traite au 2e passage que les 18\,137 nouveaux messages, et trouve pourtant 1\,919
incidents, exactement comme le batch qui relit les 70\,044 messages du topic.
\end{resultat}

\begin{obtenir}
  \item Dans \nb{05\_batch\_ou\_streaming}, exécutez toutes les cellules, dans l'ordre.
  \item Capturez la sortie de la cellule 9, le tableau du bilan.
\end{obtenir}
\capture{tp4_05_bilan}{Partie 5 : le bilan des trois façons de calculer les incidents par transporteur.}{fig:res05}
\begin{lecture}
Le flux continu rend l'incident visible 5,1 secondes après son envoi, avec des micro-lots d'environ
1 seconde, mais il a créé 5 snapshots en moins d'une minute et garde ses ressources réservées jour
et nuit. Le batch relit les 70\,044 messages en 1,1 seconde, et le flux déclenché traite les seuls
nouveaux messages en 2,1 secondes. À ce volume, relire tout ne coûte presque rien. Mais la durée du
batch grandit avec l'historique, alors que celle du flux déclenché ne dépend que des nouveaux
messages. Le choix se fait donc sur le délai dont le métier a besoin : quelques secondes pour
prévenir un client, une heure pour un tableau de bord, un mois pour un bilan.
\end{lecture}


% =====================================================================
\section{En entreprise}
% =====================================================================

\begin{center}
\small
\renewcommand{\arraystretch}{1.25}
\begin{tabularx}{\textwidth}{@{}l X X@{}}
\toprule
& \textbf{Dans ce TP} & \textbf{En production} \\
\midrule
Brokers & un seul broker, sans réplication & au moins trois brokers, réplication de 3 (Confluent Cloud, Amazon MSK, Redpanda) \\
Format des messages & JSON, schéma écrit dans le notebook & Avro ou Protobuf, schémas versionnés dans un \emph{Schema Registry} \\
Base de données vers Kafka & rejeu de fichiers & capture des changements (CDC) avec Debezium \\
Traitement & Spark Structured Streaming & Spark, Flink, Kafka Streams \\
Surveillance & Kafka UI & retard des groupes et débit dans Grafana, alertes \\
\bottomrule
\end{tabularx}
\end{center}

\needspace{24\baselineskip}
% =====================================================================
\section{Livrables}\label{sec:livrables}
% =====================================================================

\subsection*{Le compte rendu}

Au format PDF, nommé \cmd{NOM\_Prenom\_TP4.pdf}, il ne recopie pas le code mais explique vos choix
et vos résultats. Il contient :
\begin{enumerate}[leftmargin=*,itemsep=2pt]
  \item \textbf{Architecture} : votre schéma du flux, de la source aux usages, et le rôle de chaque composant.
  \item \textbf{Kafka} : le rôle des partitions, de la clé et des groupes, avec les captures de Kafka UI.
  \item \textbf{Garanties} : ce qui se passe pendant et après la panne, et le rôle du checkpoint.
  \item \textbf{Le temps} : le désordre mesuré, votre choix de watermark et son effet sur les
        incidents par jour.
  \item \textbf{Batch ou streaming} : votre tableau de mesures, et votre choix pour chacun des quatre
        besoins de la question~3 du notebook 05.
  \item \textbf{Réponses aux questions} des notebooks, rédigées et argumentées.
  \item \textbf{Limites et améliorations} : ce qu'il faudrait de plus pour passer en production
        (réplication, schémas, surveillance du retard, coûts).
\end{enumerate}

\subsection*{Les notebooks}

Les cinq notebooks de \cmd{notebooks/mon\_travail/tp4/}, exécutés, avec leurs sorties.

\clearpage
% =====================================================================
\section*{Annexe -- Aide-mémoire}\label{sec:aidememoire}
% =====================================================================

\begin{center}
\small
\renewcommand{\arraystretch}{1.22}
\begin{tabularx}{\textwidth}{@{}>{\raggedright\arraybackslash}p{4.1cm} >{\raggedright\arraybackslash\ttfamily\footnotesize}X@{}}
\toprule
\normalfont\textbf{Je veux\dots} & \normalfont\textbf{Code} \\
\midrule
\multicolumn{2}{@{}l}{\textbf{\textcolor{lhblue}{Services et rejeu}}} \\
Démarrer Kafka et Kafka UI & docker compose -{}-profile streaming up -d kafka kafka-ui \\
Rejouer un mois dans Kafka & python /opt/lakehouse/scripts/rejouer\_suivi.py -{}-depuis 2018-03-01 -{}-jusqu-a 2018-04-01 \\
Plus vite, ou d'un coup & \dots{} -{}-jours-par-seconde 5 \quad \dots{} -{}-sans-pause \\
Dans un dossier & \dots{} -{}-dossier s3://landing/flux/suivi \\
\midrule
\multicolumn{2}{@{}l}{\textbf{\textcolor{lhblue}{Kafka en Python}}} \\
Créer un topic & admin.create\_topics([NewTopic("t", num\_partitions=3, replication\_factor=1)]) \\
Lire des messages & c = Consumer(\{"bootstrap.servers": KAFKA, "group.id": "g"\}) \quad c.subscribe(["t"]) \quad c.poll(5) \\
Valider la position & c.commit(asynchronous=False) \\
Fin d'une partition & c.get\_watermark\_offsets(TopicPartition("t", 0)) \\
\midrule
\multicolumn{2}{@{}l}{\textbf{\textcolor{lhblue}{Spark Structured Streaming}}} \\
Lire Kafka & spark.readStream.format("kafka").option("kafka.bootstrap.servers", KAFKA).option("subscribe", "t").option("startingOffsets", "earliest").load() \\
Lire une table Iceberg & spark.readStream.table("lakehouse.bronze.suivi\_evenements") \\
Décoder le JSON & F.from\_json(F.col("value").cast("string"), SCHEMA) \\
Écrire dans Iceberg & df.writeStream.option("checkpointLocation", chemin).toTable("lakehouse.bronze.x") \\
Déclencheur & .trigger(processingTime="5 seconds") \quad .trigger(availableNow=True) \\
Mode de sortie & .outputMode("append") \quad "update" \quad "complete" \\
Attendre les retardataires & .withWatermark("horodatage\_ts", "12 hours") \\
Fenêtre de temps & .groupBy(F.window("horodatage\_ts", "1 day"), "transporteur").count() \\
Écarter les doublons & .dropDuplicatesWithinWatermark(["evenement\_id"]) \\
Mettre à jour une table & .writeStream.foreachBatch(fonction) \textrm{puis} MERGE INTO \dots \\
Suivre une requête & q.recentProgress \quad q.processAllAvailable() \quad q.stop() \\
\midrule
\multicolumn{2}{@{}l}{\textbf{\textcolor{lhblue}{Iceberg}}} \\
Snapshots, fichiers & SELECT * FROM lakehouse.bronze.suivi\_evenements.snapshots \quad \dots .files \\
Purger l'historique & CALL lakehouse.system.expire\_snapshots(table => '...', older\_than => TIMESTAMP '...', retain\_last => 1) \\
\bottomrule
\end{tabularx}
\end{center}

\end{document}
